\chapter{Türbülans ve Large Eddy Simulation}
\label{ch:turbulans_les}

Türbülans, ak\i{}skanlar dinamiğinin en derin ve çözülmemiş problemidir.
Bu bölümde türbülansin fizigini, Kolmogorov teorisini, DNS'in sinirlarini ve bu sinirlari aşmak için geliştirilen LES yöntemini inceliyoruz --- çünkü INNATE modelimiz bir LES çözücüsü olarak calisacak.


% ============================================================
\section{Türbülans Nedir?}
\label{sec:turbulans_nedir}

\subsection{Temel Özellikler}

Türbülans, yüksek Reynolds sayisinda ortaya çikan \textbf{düzensiz, kaotik, çok-ölçekli} ak\i{}s rejimidir. Tanimlayici özellikleri:

\begin{enumerate}
\item \textbf{Düzensizlik (irregularity):} Hiz, basinc ve sicaklik alanlari hem zamanda hem uzayda düzensiz dalgalanma gösterir. Tek bir realizasyon tahmin edilemez, ama istatistiksel ortalamalar anlamlidir.

\item \textbf{Çok-ölçeklilik (multi-scale):} Enerji taşiyan büyük girdaplardan ($\sim L$, integral ölçek) enerji yok eden en küçük girdaplara ($\sim \eta$, Kolmogorov ölçegi) kadar geniş bir ölçek araliği vardir.

\item \textbf{3B girdap gerilmesi (vortex stretching):} 2B türbülanstan temelden farkli bir mekanizma. Girdap çizgileri gerilir, incelir ve yoğunlaşir:
\begin{equation}
\DDt{\vect{\omega}} = (\vect{\omega}\cdot\gradient)\vect{u} + \nu\laplacian\vect{\omega}
\label{eq:vorticity_eq}
\end{equation}
$(\vect{\omega}\cdot\gradient)\vect{u}$ terimi ``girdap gerilmesi''dir --- sadece 3B'de vardir (2B'de $\omega$ skalerdir, bu terim s\i{}fir olur).

\item \textbf{Dissipasyon:} Türbülans sürekli kinetik enerjiyi isiya dönüştürür. Enerji beslemezseniz (forcing), türbülans zamanla söner (Taylor-Green vortex decay problemi tam olarak budur).

\item \textbf{Difuzivite artisi:} Türbülans, momentum ve isi transferini moleküler difuzyondan çok daha etkili yapar. Bu yüzden türbülansl\i{} akislarda Nusselt sayisi yüksektir.
\end{enumerate}

\subsection{Laminar--Türbülant Geçiş}

Reynolds sayısı arttıkça akış rejimi değişir:
\begin{itemize}
\item $\Rey < \Rey_{\text{cr}}$: Laminar akis (düzenli, tahmin edilebilir)
\item $\Rey \approx \Rey_{\text{cr}}$: Geçis (instabilite, Kelvin-Helmholtz girdaplari)
\item $\Rey \gg \Rey_{\text{cr}}$: Tam gelismis türbülans
\end{itemize}

Boru akisi için $\Rey_{\text{cr}} \approx 2300$, düz levha için $\Rey_{\text{cr}} \approx 5\times10^5$.
Bizim periyodik kutumuzda kritik $\Rey$ geometriye ve forcing'e bağli: $\Rey \approx 1000$'de türbülansi görmeye baslariz.


% ============================================================
\section{Kolmogorov Teorisi (1941)}
\label{sec:kolmogorov}

Andrey Kolmogorov'un 1941'deki çalişmasi (K41), türbülans teorisinin temel taşidir.

\subsection{Temel Varsayimlar}

\begin{enumerate}
\item \textbf{Yeterince yüksek $\Rey$'de} (küçük ölçeklerde), türbülans \textbf{istatistiksel olarak izotropik}tir. Büyük ölçeklerdeki anizotropi (mesela forcing yönü) küçük ölçeklere aktarilirken ``unutulur''.
\item Küçük ölçeklerin istatistikleri \textbf{evrensel}dir: sadece iki parametreye bağlidir:
\begin{itemize}
\item $\varepsilon$: enerji dissipasyon hizi $[\text{m}^2/\text{s}^3]$
\item $\nu$: kinematik viskozite $[\text{m}^2/\text{s}]$
\end{itemize}
\end{enumerate}

\subsection{Enerji Spektrumu: $-5/3$ Yasasi}
\label{subsec:energy_spectrum}

Enerji spektrumu $E(k)$, dalga sayisi $k$'deki birim dalga sayisi basina kinetik enerji yoğunluğudur:
\begin{equation}
\frac{1}{2}\langle |\vect{u}|^2 \rangle = \int_0^\infty E(k)\,dk
\label{eq:total_energy}
\end{equation}

\textbf{Inertial subrange}'de (enerji üretiminden ve dissipasyondan uzak bölge), $E(k)$ sadece $\varepsilon$ ve $k$'ya bağlıdır:
\begin{equation}
E(k) = f(\varepsilon, k) \qquad\text{(fonksiyonel form bilinmiyor, boyut analizi yapalim)}
\end{equation}

\begin{proof}[Türetme --- Boyut Analizi]
Birimleri yazalim:
\begin{align}
[E(k)] &= \frac{\text{enerji}}{\text{dalga sayisi}} = \frac{\text{m}^2/\text{s}^2}{1/\text{m}} = \text{m}^3/\text{s}^2 \label{eq:dim_Ek} \\
[\varepsilon] &= \text{m}^2/\text{s}^3 \label{eq:dim_eps} \\
[k] &= 1/\text{m} \label{eq:dim_k}
\end{align}

$E(k) = C\,\varepsilon^\alpha\, k^\beta$ formunda yazalim. Birimleri eslestir:
\begin{equation}
\text{m}^3\,\text{s}^{-2} = (\text{m}^2\,\text{s}^{-3})^\alpha \cdot (\text{m}^{-1})^\beta
\end{equation}

Metre üs: $3 = 2\alpha - \beta$. Saniye üs: $-2 = -3\alpha$, dolayisiyla $\alpha = 2/3$.

$\beta = 2\alpha - 3 = 4/3 - 3 = -5/3$. Sonuç:

\begin{equation}
\boxed{E(k) = C_K\,\varepsilon^{2/3}\,k^{-5/3}}
\label{eq:kolmogorov_spectrum}
\end{equation}

$C_K \approx 1.5$ \textbf{Kolmogorov sabiti}dir (deneysel olarak belirlenmistir).
\end{proof}

Bu iki-üçlü yasasi, türbülans teorisinin en önemli ve en çok dogrulanan sonuçudur. Atmosferden okyanuslara, y\i{}ldiz içlerinden laboratuvar deneylerine kadar geçerlidir.

\begin{figure}[H]
\centering
\includegraphics[width=0.85\textwidth]{energy_spectrum.png}
\caption{Türbülant enerji spektrumu $E(k)$. Üç bölge: enerji enjeksiyonu (düşük $k$), inertial subrange ($E(k) \sim k^{-5/3}$), ve dissipasyon (yüksek $k$). DNS tüm bölgeleri çözer; LES yalnızca kesme çizigisinin solunu çözer.}
\label{fig:energy_spectrum}
\end{figure}

% ============================================================
\section{Enerji Kaskadi}
\label{sec:cascade}

\subsection{Kavramsal Resim}

Türbülanstaki enerji ak\i{}s\i{} tek yönlüdür (3B'de):

\begin{equation}
\underbrace{\text{Enerji üretimi}}_{\text{büyük ölçek } (k \sim k_f)}
\quad\xrightarrow{\text{kaskad}}\quad
\underbrace{\text{Inertial transfer}}_{\text{orta ölçekler}}
\quad\xrightarrow{\text{kaskad}}\quad
\underbrace{\text{Dissipasyon}}_{\text{küçük ölçek } (k \sim k_\eta)}
\end{equation}

Detayli olarak üç bölge vardir:

\begin{enumerate}
\item \textbf{Energy-containing range} ($k \sim k_f$): D\i{}s forcing veya mean shear ile enerji enjekte edilir. Büyük girdaplar burada yasarlar. $E(k)$ davranisi forcing'e bağli (evrensel degil).

\item \textbf{Inertial subrange} ($k_f \ll k \ll k_\eta$): Enerji ne üretilir ne yok edilir, sadece büyük girdaplardan küçüklere \textbf{transfer} edilir. Kolmogorov $-5/3$ yasasi tam burada geçerlidir. Enerji transfer hizi sabittir: $\Pi(k) = \varepsilon$.

\item \textbf{Dissipation range} ($k \gtrsim k_\eta$): Viskoz kuvvetler baskindir. Kinetik enerji isiya donusur:
\begin{equation}
\varepsilon = 2\nu\,\mathcal{Z}, \qquad \mathcal{Z} = \frac{1}{2}\langle|\vect{\omega}|^2\rangle \;\text{(enstrofi)}
\label{eq:dissipation_enstrophy}
\end{equation}
\end{enumerate}

\begin{figure}[H]
\centering
\includegraphics[width=0.8\textwidth]{energy_cascade.png}
\caption{Enerji kaskadı şematik gösterimi. Soldan sağa: enerji üretimi, inertial transfer ($-5/3$ yasası), ve dissipasyon.}
\label{fig:energy_cascade}
\end{figure}


% ============================================================
\section{Kolmogorov Uzunluk Ölçegi}
\label{sec:kolmogorov_scale}

En küçük girdabin boyutu ne kadardir? Dissipasyon bölgesinde iki mekanizma yarışır: inertial transfer ($\varepsilon$) ve visköz difüzyon ($\nu$). Boyut analizi ile:

\begin{equation}
\eta = \left(\frac{\nu^3}{\varepsilon}\right)^{1/4}
\label{eq:kolmogorov_eta}
\end{equation}

\begin{proof}[Türetme]
$\eta$ bir uzunluktur $[\text{m}]$, ve sadece $\nu\;[\text{m}^2/\text{s}]$ ile $\varepsilon\;[\text{m}^2/\text{s}^3]$'a bağlidir:
\begin{align}
[\eta] &= [\nu^\alpha\,\varepsilon^\beta] \notag \\
\text{m} &= (\text{m}^2\,\text{s}^{-1})^\alpha \cdot (\text{m}^2\,\text{s}^{-3})^\beta
\end{align}
Metre: $1 = 2\alpha + 2\beta$. Saniye: $0 = -\alpha - 3\beta$, yani $\alpha = -3\beta$. Birinci denklemde: $1 = -6\beta + 2\beta = -4\beta$, dolayisiyla $\beta = -1/4$, $\alpha = 3/4$.
\begin{equation}
\eta = \nu^{3/4}\,\varepsilon^{-1/4} = \left(\frac{\nu^3}{\varepsilon}\right)^{1/4} \quad\checkmark
\end{equation}
\end{proof}

\subsection{$\eta$ ve $\Rey$ İlişkisi}

Integral ölçek $L$ ve büyük ölçek hiz $U$ ile $\varepsilon \sim U^3/L$ (boyutsal tahmin). O zaman:
\begin{equation}
\frac{\eta}{L} \sim \left(\frac{\nu^3}{U^3/L}\right)^{1/4}\cdot\frac{1}{L} = \left(\frac{\nu}{UL}\right)^{3/4} = \Rey^{-3/4}
\label{eq:eta_over_L}
\end{equation}

\begin{equation}
\boxed{\frac{\eta}{L} \sim \Rey^{-3/4}}
\end{equation}

$\Rey$ arttıkça en küçük girdap hızla küçülür: ölçek ayrımı genişler.


% ============================================================
\section{DNS (Direct Numerical Simulation)}
\label{sec:dns}

\subsection{Temel Fikir}

DNS, Navier-Stokes denklemlerini \textbf{hiçbir model kullanmadan} dogrudan çözer. Tüm ölçekler --- $\eta$'dan $L$'ye kadar --- çözünür.

\subsection{Grid Gereksinimleri}

DNS gridinin $\eta$'yi çözümleyecek kadar ince olmasi gerekir: $\Delta x \lesssim \eta$. Bu durumda her yönde gerekli nokta sayisi:
\begin{equation}
N \sim \frac{L}{\eta} \sim \Rey^{3/4}
\end{equation}

3D'de toplam nokta sayisi:
\begin{equation}
\boxed{N_{\text{total}} = N^3 \sim \Rey^{9/4}}
\label{eq:dns_cost}
\end{equation}

\begin{example}
Say\i{}sal tahminler:
\begin{center}
\begin{tabular}{ccc}
\toprule
$\Rey$ & $N_{\text{total}} \sim \Rey^{9/4}$ & Durum \\
\midrule
$1{,}000$ & $\sim 5.6\times10^6$ & Masaustu bilgisayar \\
$3{,}000$ & $\sim 8.8\times10^7$ & İstasyon/kucuk cluster \\
$10{,}000$ & $\sim 5.6\times10^9$ & Süper bilgisayar \\
$100{,}000$ & $\sim 10^{12}$ & Dünyanin en buyuk bilgisayarlari \\
\bottomrule
\end{tabular}
\end{center}

$\Rey = 10{,}000$ için: $\eta / L \sim 10{,}000^{-3/4} \approx 10^{-3}$. Eğer $L = 10$ (bizim $L_y$) ise $\eta \approx 0.01$ ve grid $N \approx 1000$. 3D'de $10^9$ nokta. Her zaman adimi için 3D FFT: $\sim 10^{10}$ islem. Binlerce zaman adimi ile çarpilinca --- birkac gün süper bilgisayarda.
\end{example}

\subsection{Bizim DNS'imiz}

Projede DNS, \textbf{validasyon} için kullanilir (INNATE'in dogru sonuc verip vermedigini kontrol etmek için). DNS'i düsük $\Rey$'de yapariz:
\begin{itemize}
\item $\Rey = 1000$: Grid $\approx 192\times320\times128$ (yapılabilir)
\item $\Rey = 3000$: Grid $\approx 384\times640\times256$ (zorlayici ama mumkun)
\item $\Rey = 10000$: Grid $\approx 2048\times3200\times1280$ --- $\sim8.4$ milyar nokta, \textbf{imkansiz} (bizim kaynaklarla)
\end{itemize}

Yüksek $\Rey$'de INNATE, DNS yerine LES yaparak sonuc uretir. DNS sadece referans olarak dusuk $\Rey$'de kullanilir.


% ============================================================
\section{LES (Large Eddy Simulation)}
\label{sec:les}

\subsection{Temel Fikir}

DNS'in $\Rey^{9/4}$ maliyeti kaldirilamaz. LES'in cozumu:
\begin{quote}
Büyük girdaplari \textbf{dogrudan çöz}, küçük girdaplari \textbf{modelle}.
\end{quote}

Fiziksel gerekce: büyük girdaplar geometriye baglidir ve tahmin edilemez; küçük girdaplar ise (Kolmogorov'a gore) evrensel ve izotropiktir, dolayisiyla basit modellerle iyi temsil edilebilir.

\subsection{Filtreleme}

LES'te hız alanı bir \textbf{uzaysal filtre} ile ayrıştırılır:
\begin{equation}
\bar{u}_i(\vect{x},t) = \int G(\vect{x}-\vect{x}',\Delta)\,u_i(\vect{x}',t)\,d\vect{x}'
\label{eq:les_filter}
\end{equation}

Burada $G$ filtre fonksiyonu, $\Delta$ filtre genisligi (grid boyutu). Filtreleme:
\begin{itemize}
\item \textbf{Cozunen (resolved)} kisim: $\bar{u}_i$ --- büyük girdaplar
\item \textbf{Cozunmeyen (subgrid-scale, SGS)} kisim: $u'_i = u_i - \bar{u}_i$ --- kucuk girdaplar
\end{itemize}

\subsection{Filtrelenmis Navier-Stokes Denklemleri}

NS denklemlerini filtreleyelim. Sureklilik denklemi filtreleme ile degismez:
\begin{equation}
\pd{\bar{u}_i}{x_i} = 0
\end{equation}

Momentum denklemi:
\begin{equation}
\pd{\bar{u}_i}{t} + \pd{(\bar{u}_i\bar{u}_j)}{x_j} = -\frac{1}{\rho}\pd{\bar{p}}{x_i} + \nu\pdd{\bar{u}_i}{x_j} - \pd{\tau_{ij}^{\text{sgs}}}{x_j}
\label{eq:les_ns}
\end{equation}

Burada $\tau_{ij}^{\text{sgs}}$ \textbf{sub-grid scale (SGS) stres tensorudur}:
\begin{equation}
\tau_{ij}^{\text{sgs}} = \overline{u_i u_j} - \bar{u}_i\bar{u}_j
\label{eq:sgs_stress}
\end{equation}

Bu terim \textbf{bilinmiyor} --- cozunmeyen küçük olceklerin cozunen olceklere etkisini temsil eder. LES'in tum zorluğu bu terimi modellemektir.

\begin{remark}
SGS stresinin fiziksel anlami: küçük girdaplarin buyuk girdaplara uyguladiği ``sürtünme kuvveti''. Net etki: enerjiyi küçük olceklere transfer etmek (dissipasyon sağlamak).
\end{remark}


% ============================================================
\section{Smagorinsky Modeli}
\label{sec:smagorinsky}

En basit ve en yaygin SGS modeli, Smagorinsky (1963) tarafindan atmosfer modellemesi için önerilmiştir.

\subsection{Eddy Viscosity Hipotezi}

SGS stresinin, Boussinesq eddy-viscosity hipotezi ile modellenebileceği varsayilir:
\begin{equation}
\tau_{ij}^{\text{sgs}} - \frac{1}{3}\tau_{kk}^{\text{sgs}}\,\delta_{ij} = -2\nu_t\,\bar{S}_{ij}
\label{eq:eddy_viscosity}
\end{equation}

Burada $\nu_t$ \textbf{türbülans (eddy) viskozitesidir} --- molekuler viskoziteye ek olarak, küçük girdaplarin etkisini temsil eder. $\bar{S}_{ij}$ filtrelenmis strain rate tensoru:
\begin{equation}
\bar{S}_{ij} = \frac{1}{2}\left(\pd{\bar{u}_i}{x_j} + \pd{\bar{u}_j}{x_i}\right)
\label{eq:strain_rate}
\end{equation}

\subsection{Smagorinsky Formulasyonu}

$\nu_t$ boyut analizi ile belirlenir:
\begin{equation}
[\nu_t] = \text{m}^2/\text{s} = [\text{uzunluk}]^2 \cdot [\text{zaman}]^{-1}
\end{equation}

Elimizdeki uzunluk olcegi: filtre genisligi $\Delta$. Zaman olcegi: $1/|\bar{S}|$. Dolayisiyla:
\begin{equation}
\boxed{\nu_t = (C_s\,\Delta)^2\,|\bar{S}|}
\label{eq:smagorinsky}
\end{equation}

Burada:
\begin{itemize}
\item $C_s \approx 0.1\text{--}0.2$: \textbf{Smagorinsky sabiti} (boyutsuz)
\item $\Delta = (\Delta x\cdot\Delta y\cdot\Delta z)^{1/3}$: filtre genisligi (kup kok)
\item $|\bar{S}| = \sqrt{2\,\bar{S}_{ij}\bar{S}_{ij}}$: strain rate magnitude
\end{itemize}

\begin{remark}
$C_s$'nin degeri tartismalidir. Homojen izotropik turbulans icin $C_s \approx 0.17$ (Lilly, 1967'nin teorik tahmini). Kanal akisi icin $C_s \approx 0.1$ daha uygun. INNATE'te $C_s$ \textbf{ogrenilebilir bir parametredir} --- model, probleme ozgu optimal degeri kendi bulur.
\end{remark}

\subsection{Efektif Viskozite}

LES'te toplam viskoz etki:
\begin{equation}
\nu_{\text{eff}} = \nu + \nu_t
\label{eq:nu_eff}
\end{equation}

Filtrelenmis momentum denklemi:
\begin{equation}
\pd{\bar{u}_i}{t} + (\bar{\vect{u}}\cdot\gradient)\bar{u}_i = -\frac{1}{\rho}\pd{\bar{p}}{x_i} + \nu_{\text{eff}}\laplacian\bar{u}_i
\label{eq:les_effective}
\end{equation}

Bu form, DNS denkleminin \emph{aynısı}'dır --- tek fark $\nu$ yerine $\nu_{\text{eff}}$ kullanılması. Çok zarif: LES, DNS kodunun üzerine sadece $\nu_t$ hesabi eklenerek uygulanabilir.

\subsection{Strain Rate Magnitude Hesabi}

$|\bar{S}|$'yi acik olarak yazmak gerekirse (3D'de):

\begin{align}
\bar{S}_{ij}\bar{S}_{ij} &= \bar{S}_{xx}^2 + \bar{S}_{yy}^2 + \bar{S}_{zz}^2 + 2(\bar{S}_{xy}^2 + \bar{S}_{xz}^2 + \bar{S}_{yz}^2) \label{eq:sij_sij}
\end{align}

Her bilesen:
\begin{align}
\bar{S}_{xx} &= \pd{\bar{u}}{x}, \qquad \bar{S}_{yy} = \pd{\bar{v}}{y}, \qquad \bar{S}_{zz} = \pd{\bar{w}}{z} \\
\bar{S}_{xy} &= \frac{1}{2}\left(\pd{\bar{u}}{y}+\pd{\bar{v}}{x}\right), \quad
\bar{S}_{xz} = \frac{1}{2}\left(\pd{\bar{u}}{z}+\pd{\bar{w}}{x}\right), \quad
\bar{S}_{yz} = \frac{1}{2}\left(\pd{\bar{v}}{z}+\pd{\bar{w}}{y}\right)
\end{align}

\begin{equation}
|\bar{S}| = \sqrt{2\,\bar{S}_{ij}\bar{S}_{ij}}
\label{eq:strain_mag}
\end{equation}

Tüm bu türevler spectral yöntemle (Bölüm~\ref{ch:spectral}) hesaplanir.


% ============================================================
\section{INNATE'te EddyViscosity3D}
\label{sec:eddy_viscosity_innate}

Projemizdeki \texttt{EddyViscosity3D} noronu, Smagorinsky modelinin INNATE implementasyonudur.

\subsection{Noron Yapisi}

\begin{equation}
\nu_t = \big(\underbrace{\tilde{C}_s}_{\text{ogrenilebilir}}\cdot\Delta\big)^2 \cdot |\bar{S}|
\end{equation}

\textbf{Tek ogrenilebilir parametre:} $\tilde{C}_s$ (başlangıç değeri $0.15$, $[0.05, 0.30]$ arasinda clamp edilir).

\textbf{Grid spacing:}
\begin{equation}
\Delta = (\Delta x \cdot \Delta y \cdot \Delta z)^{1/3} = \left(\frac{L_x}{N_x}\cdot\frac{L_y}{N_y}\cdot\frac{L_z}{N_z}\right)^{1/3}
\end{equation}

Bizim grid icin: $\Delta = (6/96 \cdot 10/160 \cdot 4/64)^{1/3} = (0.0625 \cdot 0.0625 \cdot 0.0625)^{1/3} = 0.0625$.

\subsection{Neden LES Zorunlu?}

63 saatlik basarisiz egitimin \textbf{birincil sebebi} EddyViscosity3D'nin KAPALI olmasiydi. Grid analizi sonuclari bunu dogrular:

\begin{center}
\begin{tabular}{lcc}
\toprule
Grid & $\Rey = 5{,}000$ & $\Rey = 10{,}000$ \\
\midrule
$96\times160\times64$ (LES skoru) & 80/100 (iyi) & 74/100 (makul) \\
$32\times48\times24$ (LES skoru) & 60/100 (zayif) & 44/100 (yetersiz) \\
DNS gereksinimleri & $\sim10^8$ nokta & $\sim8.4\times10^9$ nokta \\
\bottomrule
\end{tabular}
\end{center}

$96\times160\times64 = 983{,}040$ nokta. $\Rey = 10{,}000$ için DNS $\sim 8.4\times10^9$ nokta gerektirir --- $8{,}500\times$ daha fazla. LES olmadan bu gridde yüksek $\Rey$ çözülemez: subgrid-scale dissipasyon saglanamaz ve enerji birikir $\to$ sayisal patlama (NaN).


% ============================================================
\section{LES Ötesinde: Gelismis SGS Modelleri}
\label{sec:advanced_sgs}

Smagorinsky modeli basit ve saglamdir ama sinirliliklari vardir:

\begin{itemize}
\item Duvar yakininda asiri dissipasyon verir ($C_s$'yi küçültmek gerekir).
\item Laminar--türbülant geçiste türbülans geçiş başlatamaz (``backscatter'' yoktur).
\item $C_s$ sabittir --- aslinda akisa ve konuma baglidir.
\end{itemize}

Gelismis modeller:
\begin{description}
\item[Dynamic Smagorinsky (Germano, 1991):] $C_s$'yi lokal olarak akinstan hesaplar (test filtre kullanarak). Cok basarili ama hesaplama maliyeti yüksek.
\item[WALE (Nicoud \& Ducros, 1999):] Duvar yakininda dogru asimptotik davranis. Ek turev hesabi gerektirir.
\item[Sigma model (Nicoud et al., 2011):] Hiz gradient tensorunun tekil degerleri ile $\nu_t$ hesaplar. Duvar yakininda ve laminar akista dogru sonuc verir.
\end{description}

\begin{remark}
INNATE'te $C_s$ zaten ogrenilebilir oldugundan, modelin Dynamic Smagorinsky'nin isini ``ogrenme yoluyla'' yapmasini bekliyoruz. Lokal-adaptif parametreler kullanilirsa (Yol~C mimarisi), $C_s(\vect{x},t)$ uzaysal olarak degisir --- bu, Dynamic Smagorinsky'ye fonksiyonel olarak esittir.
\end{remark}


% ============================================================
\section{Bu Bölümün Özeti}

\begin{enumerate}
\item \textbf{Türbülans}: Duzensiz, cok-olcekli, kaotik akis. $\Rey$ arttikca ortaya cikar.
\item \textbf{Kolmogorov teorisi}: Inertial subrange'de $E(k) = C_K\varepsilon^{2/3}k^{-5/3}$ (evrensel).
\item \textbf{Kolmogorov olcegi}: $\eta = (\nu^3/\varepsilon)^{1/4}$, $\eta/L \sim \Rey^{-3/4}$.
\item \textbf{DNS maliyeti}: $N^3 \sim \Rey^{9/4}$ --- yüksek $\Rey$'de imkansiz.
\item \textbf{LES}: Buyuk girdaplari coz, kucukleri modelle. SGS stresi: $\tau_{ij}^{\text{sgs}} = \overline{u_iu_j} - \bar{u}_i\bar{u}_j$.
\item \textbf{Smagorinsky}: $\nu_t = (C_s\Delta)^2|\bar{S}|$ --- basit, sağlam, ama sabit $C_s$.
\item \textbf{INNATE'te LES}: EddyViscosity3D noronu, ogrenilebilir $C_s$ ile Smagorinsky modeli uygular. Grid analizi: $96\times160\times64$ gridde $\Rey = 10{,}000$ için LES zorunlu.
\end{enumerate}

Bir sonraki bolumde, turbulans tahmini icin gelistirilen neural operator yaklasimlarini --- PINN, FNO, DeepONet ve bizim INNATE konseptimizi --- inceliyoruz.
